\documentclass[9pt,a4paper,twoside]{rho-class/rho}
%%%%%%%%%%%%%%%%%%%%%%%%%%%%%%%%%%%%%%%%%%%%%%%%%%%%%%%%%%%
% --------------------------------------------------------
% Rho
% LaTeX Template
% Version 2.1.1 (01/09/2024)
%
% Authors: 
% Guillermo Jimenez (memo.notess1@gmail.com)
% Eduardo Gracidas (eduardo.gracidas29@gmail.com)
% 
% License:
% Creative Commons CC BY 4.0
% --------------------------------------------------------
%%%%%%%%%%%%%%%%%%%%%%%%%%%%%%%%%%%%%%%%%%%%%%%%%%%%%%%%%%%

%----------------------------------------------------------
% Font Config
%----------------------------------------------------------
% 加载 ctex (用于翻译和中文排版)
% 1. 引入 ctex 宏包（它会自动处理 xeCJK，不需要再次 \usepackage{xeCJK}）
\usepackage[UTF8]{ctex}
\usepackage{indentfirst}

% 2. 引入 unicode-math 用于处理数学字体
% 注意：请在 setmainfont 之前加载
\usepackage{unicode-math}

% 3. 设置西文字体 (Times New Roman)
\setmainfont{Times New Roman}
\setsansfont{Arial}
\setmonofont{Courier New}

% 4. 设置数学字体 (关键步骤！)
% TeX Gyre Termes Math 是 Times New Roman 风格的数学字体，Overleaf 内置支持
\setmathfont{TeX Gyre Termes Math}

% 5. 设置中文字体
% 注意：Overleaf 的 Noto 字体路径可能需要根据文件名精确指定，或者使用 Fandol 系列（ctex默认）
% 这里保留你的设置，但在本地编译时可能需要调整
\setCJKmainfont[AutoFakeBold=true]{Noto Serif CJK SC}
\setCJKsansfont[AutoFakeBold=true]{Noto Sans CJK SC}
\setCJKmonofont{Noto Serif CJK SC} % 建议等宽字体改为 Noto Sans Mono CJK SC

\usepackage{float}
\usepackage{array}
\usepackage{booktabs}
\usepackage{caption}
\usepackage{tabularx}

% 1.5倍行距
\setstretch{1.5}

% 设置专属于 PBL 的病历简述
\renewcommand{\abstractname}{病例摘要}
\renewcommand{\keywordname}{\textbf{关键词}：}
\newcommand{\namebox}[1]{\makebox[3em][s]{#1}}

%% Spanish babel recomendation
% \usepackage[spanish,es-nodecimaldot,es-noindentfirst]{babel}

\setbool{rho-abstract}{true} % Set false to hide the abstract
\setbool{corres-info}{false} % Set false to hide the corresponding author section

%----------------------------------------------------------
% TITLE
%----------------------------------------------------------

\journalname{PKU School of Basic Medicine Science}
\title{PBL-G1：甲型流感}

%----------------------------------------------------------
% AUTHORS AND AFFILIATIONS
%----------------------------------------------------------

\author[1]{文康丞}

%----------------------------------------------------------

\affil[1]{24级临床五班，北京大学医学部}

%----------------------------------------------------------
% DATES
%----------------------------------------------------------

\dates{本报告文档于 2025年12月7日 编译}

%----------------------------------------------------------
% FOOTER INFORMATION
%----------------------------------------------------------

\leadauthor{Syncrex Wen}
\footinfo{PBL-G1}
\smalltitle{流行性感冒}
\institution{Peking University}
\theday{Dec 7, 2025} %\today

%----------------------------------------------------------
% ARTICLE INFORMATION
%----------------------------------------------------------
% SYNCREX: I don't need this section so I just comment it
% \corres{Provide the corresponding author information and publisher here.}
% \email{example@organization.com.}
% \doi{\url{https://www.doi.org/exampledoi/XXXXXXXXXX}}

% \received{March 20, 2024}
% \revised{April 16, 2024}
% \accepted{April 20, 2024}
% \published{May 21, 2024}

% 这是原模版的版权信息
\license{Rho LaTeX Class \ccLogo\ This document is licensed under Creative Commons CC BY 4.0. Modified by Syncrex Wen to support Chinese.}

%----------------------------------------------------------
% ABSTRACT
%----------------------------------------------------------

\begin{abstract}

    患者张某，男，65 岁，因"流涕、干咳 1 天，高热、寒战、头痛及肌肉酸痛半天"就诊。症状起初表现为流涕及刺激性干咳，随后出现高热（体温约 39$^o\mathrm{C}$） 、畏寒寒战、头痛及全身肌痛，无咳脓痰、胸痛、呼吸困难及消化道症状。既往有 30 余年吸烟史，确诊慢性阻塞性肺疾病（COPD）但未规律治疗。就诊前曾接触咳嗽患者，家中外孙所在班级因" 流感" 停课。体格检查所示咽部充血，扁桃体肿大；双肺呼吸音粗，未闻及明显干湿啰音；胸腹叩诊无异常。\\
    
    辅助检查显示甲型流感抗原阳性，新冠检测阴性；血常规示淋巴细胞下降提示病毒感染；胸部 CT 可见肺气肿及双肺少许间质性改变。综合病史、体征及检查结果，诊断为" 甲型流感病毒感染"，予奥司他韦及布洛芬治疗并指导居家隔离 5 天。复诊时患者症状明显改善，甲型流感病毒核酸检测转阴。

\end{abstract}

%----------------------------------------------------------

\keywords{流行性感冒，甲型流感}

%----------------------------------------------------------

\begin{document}

%----------------------------------------------------------
% Render the first page
%----------------------------------------------------------

    % 重置页码标签
    \pagenumbering{arabic}
    \maketitle
    \thispagestyle{firststyle}
    \tableofcontents
    \linenumbers

%----------------------------------------------------------
\section{病毒学与呼吸道感染综述}

    \subsection{病毒学}

    \subsubsection{病毒的形态与基本结构}

    病毒体大小差别悬殊，最大的长度可达 1µm 以上，最小的病毒仅十几纳米，一般介于 20-300nm 之间，大多数病毒小于 150nm，流感病毒的直径即大约 100nm。球形病毒的大小用其直径表示，其他形状病毒则以长度 ×宽度等表示。

    我们大致将病毒体的形态分为 5 类：球形(spheroid)、丝状(filament)、弹状(bullet shape)、砖状(brick shape)和蝌蚪状(tadpole shape)。大多数感染人和动物的病毒都属于球形，如脊髓灰质炎病毒、冠状病毒、人类免疫缺陷病毒、流感病毒；此外，丝状病毒科的埃博拉病毒属于丝状、初次分离的流感病毒也可呈丝状；狂犬病病毒和水疱性口炎病毒属于弹状；天花病毒和痘苗病毒属于砖状；大多数噬菌体外形呈蝌蚪状，另外，有些病毒可具有多形性。
    
    病毒体的基本结构是由核心（core）和衣壳（capsid）构成的核衣壳（nucleocapsid）。有些病毒的核衣壳外有包膜（envelope）和包膜的构成成分刺突（spike）。有包膜的病毒称为包膜病毒，无包膜的病毒体称裸露病毒，核衣壳是裸露病毒完整的病毒体。

\begin{enumerate}
    \item \textbf{核心}：位于病毒体的最内部，主要化学成分为核酸，由一种类型的核酸（DNA或RNA）组成，构成病毒基因组。此外，有些病毒体的核心含有少量蛋白质，多为携带的酶类。核心是病毒执行生命活性的物质基础。
    \item \textbf{衣壳}：由病毒基因组编码的包围在病毒核心外面的蛋白质外壳。衣壳系由一定数量的壳粒（capsomere）组成，每个壳粒被称为形态亚单位（morphologic subunit），由一个或多个多肽分子组成。病毒衣壳结构遵循对称性规律，根据所含壳粒数目和排列方式不同，病毒衣壳可分为三种不同对称型。
    \item \textbf{包膜}：部分病毒在核衣壳外围绕一层镶嵌有多糖蛋白的脂质双层膜结构，称为病毒的包膜。包膜是在病毒增殖成熟时，穿过宿主细胞（器）膜出芽释放时获得的，脂质双层和膜上的多糖分子成分都来源于宿主细胞，但包膜蛋白是由病毒基因编码的，多糖和蛋白构成包膜糖蛋白，主要作用是保护核衣壳，构成病毒的表面抗原，参与机体应答免疫过程。此外，有的病毒上还有基质蛋白和触须结构。
\end{enumerate}

    \subsubsection{病毒的化学组成}

    病毒核酸位于病毒体核心，只含有一种核酸（DNA 或 RNA），构成病毒体基因组，携带病毒遗传需要的所有信息，大小从 3.2kb 到 32.0kb 不等，携带病毒的所有遗传信息；病毒的基因组核酸类型多种多样。

    \begin{table}[H]
        \centering
        \caption{病毒基因组特性分类}
        \label{tab:virus_genome}
        \begin{tabular}{@{}llllll@{}}
        \toprule
        \textbf{特性} & \textbf{基因组/类型} & \textbf{DNA 病毒} & \textbf{举例} & \textbf{RNA 病毒} & \textbf{举例} \\
        \midrule
        形状 & 线状单链 & +ssDNA & 细小病毒 B19 & +ssRNA & 小 RNA 病毒 \\
            & 线状双链 & dsDNA & 腺病毒 & dsRNA & 呼肠病毒 \\
            & 环状单链 & +scDNA & M13 噬菌体 & / & / \\
            & 环状单链 & -scDNA & TT 病毒 & -ssRNA & 丁型肝炎病毒 \\
            & 环状双链 & dcDNA & 乳头瘤病毒 & / & / \\
        \midrule
        完整性 & 不分段 & 有 & 多数 DNA 病毒 & 有 & 多数 RNA 病毒 \\
               & 节段 & 有 & 双生病毒 & 有 & 沙粒、布尼亚、正黏及呼肠病毒 \\
        \midrule
        构成 & 单倍体 & 都是 & 腺病毒等 & 有 & RNA 病毒（除外逆转录病毒） \\
             & 双倍体 & 无 & 无 & 有 & 逆转录病毒 \\
        \midrule
        极性 & 正链 (+) & +scDNA & M13 噬菌体 & +ssRNA & 肠道病毒、黄病毒等 \\
             & 负链 (-) & -scDNA & TT 病毒 & -ssRNA & 正黏病毒 \\
             & 双义 (±) & 有 & 腺病毒相关病毒 & 少数 & 布尼亚病毒和沙粒病毒 \\
        \bottomrule
        \end{tabular}
    \end{table}

    其中：

    \begin{itemize}
        \item 基因组极性
            \begin{itemize}
                \item 正链 (+)：指病毒基因组序列与 mRNA 一致，进入细胞后可直接作为模板翻译出蛋白质；
                \item 负链 (-)：指基因组序列与 mRNA 互补，必须先转录成正链（mRNA）后才能翻译蛋白质；
                \item 双义 (±)：指同一条核酸链上，部分区域是正链性质，部分区域是负链性质。
            \end{itemize}
        \item 基因组完整性
            \begin{itemize}
                \item 不分段：指病毒的全部遗传信息都在一条连续的核酸分子上；
                \item 分段：指病毒的基因组由多条独立的核酸片段共同组成。
            \end{itemize}
        \item 常见的结构缩写
            \begin{itemize}
                \item ss (Single-stranded) 代表单链
                \item ds (Double-stranded) 代表双链
                \item sc (Single-stranded Circular) 代表环状单链
                \item dc (Double-stranded Circular) 代表环状双链
            \end{itemize}
    \end{itemize}

    病毒蛋白约占病毒体总重的 70\%，分为结构与非结构蛋白。结构蛋白（衣壳、包膜等）构成病毒实体，负责保护核酸、维持形态、吸附宿主（病毒吸附蛋白VAP）、抗原性及血凝。非结构蛋白（如酶）主要在感染细胞内表达，辅助复制。此外，病毒含宿主来源的脂类和糖类。脂类主要存在于包膜，糖类修饰则关乎致病机制及疫苗研发。

    \subsection{呼吸道感染}

    呼吸道感染（Respiratory Tract Infections, RTIs）泛指由病毒、细菌和非典型病原体（如支原体和衣原体）引起的感染性疾病。在临床上呼吸道感染分为上呼吸道感染（Upper RIT，URIT）和下呼吸道（Lower RIT，LRIT）感染两种，二者以解剖学上的环状软骨（Cricoid Cartilage）或声带（Vocal Cords）为界限。
    
    上呼吸道感染的区域包括包括鼻腔、鼻窦、咽部（Pharynx）和喉部（Larynx，声带以上）的部分，常见的感冒（包括普通感冒和流行性感冒）、鼻窦炎、咽炎、喉炎、扁桃体炎等，常见的病症主要局限于头颈部，包括鼻塞、流涕、打喷嚏、咽痛、声音嘶哑。通常伴有低热，在流感中可能出现高热。下呼吸道的区域包含解剖范围内气管（Trachea）、支气管（Bronchi）、细支气管（Bronchioles）和肺泡（Alveoli）等，下呼吸道感染常见的疾病有气管炎、支气管炎、肺炎（Pneumonia）等。下呼吸道感染的症状主要集中在胸部，包括咳嗽（干咳或伴有咳痰）、呼吸困难（气促）、胸痛、喘息，全身症状如高热、寒战、乏力通常比URTI更严重。下呼吸道感染会直接影响气体交换功能，导致低氧血症，严重时可危及生命。
    
    引起呼吸道感染的病毒通常可以分为流感病毒（Influenza）和其他病毒两类。流感病毒具体分型为四类，即甲、乙、丙、丁四型（influenza A/B/C/D），后文将作更详尽的介绍。其他病毒包括冠状病毒 (Coronaviruses, e.g., SARS-CoV-2)、鼻病毒 (Rhinovirus)等多种类型，通常不会对人体产生很大的影响。除病毒外，还有很多细菌也能造成呼吸道感染，如肺炎链球菌 (Streptococcus pneumoniae)、流感嗜血杆菌 (Haemophilus influenzae)、金黄色葡萄球菌 (Staphylococcus aureus)等。另外，肺炎支原体 (Mycoplasma pneumoniae)、嗜肺军团菌 (Legionella pneumophila)等非典型病原体（革兰氏染色下不易着色和独特的生化特性）也会引起呼吸道感染。

\section{免疫应答综述}

\subsection{免疫应答的分类}

\subsubsection{固有免疫}

        固有免疫是机体与生俱来的防御能力，又称为先天性免疫（natural immunity）或非特异性免疫（non-specific immunity）。固有免疫是机体抵御病原体入侵和对其他危险信号产生反应的第一道防线，在危险信号刺激早期数分钟到数天内快速发挥作用。由\textbf{组织屏障}以及多种\textbf{固有免疫细胞和分子}组成。固有免疫细胞通过细胞表面或细胞内的模式识别受体识别病原微生物或机体组织损伤时释放的某些共有固定成分（也称分子模式）产生应答，保护机体免受微生物入侵或者其他因素造成的损害以及维持内环境的稳定。

        \subsubsection{适应性免疫}

        适应性免疫是机体受病原体或机体组织中的某些成分刺激后获得的防御能力，又称为获得性免疫（acquired immunity）或特异性免疫（specific immunity）。自然界中能够激活和诱导适应性免疫应答并能与免疫应答产物发生特异性结合的物质称为抗原（antigen，Ag）。在免疫应答的后期，机体适应性免疫细胞—T 淋巴细胞和 B 淋巴细胞通过特异性抗原受体识别抗原成分，从而产生抗原特异性的免疫应答来清除抗原。机体接触过某种抗原发生了适应性免疫应答之后还能产生免疫记忆，使得机体再次接触同一种抗原时能够启动更为迅速和高效的特异性免疫应答。

        \subsubsection{两种免疫应答方式的比较和联系}

        \begin{table}[H]
            \RaggedRight
            \caption{先天性免疫与获得性免疫特征比较}
            \label{tab:compare_between_two_immune_response}
                \begin{tabular}{lccp{12.2cm}}
                    \toprule
                    \textbf{特征} & \textbf{先天性免疫} & \textbf{获得性免疫} \\ 
                    \midrule
                    反应速度 & 极快（几分钟到几小时） & 较慢（初次接触需数天到数周才能达到高峰） \\
                    特异性 & 非特异性，识别的是一类病原体共有的特征 & 高度特异性 \\
                    主要细胞 & 巨噬细胞、中性粒细胞、树突状细胞、NK细胞 & T淋巴细胞（细胞免疫）、B淋巴细胞（体液免疫） \\
                    免疫记忆 & 无（再次遇到相同病毒，反应强度不变） & 有（再次遇到相同病毒，反应更快、更强） \\
                    主要分子 & 补体、干扰素（Interferon, IFN）、溶菌酶、细胞因子 & 抗体（IgG, IgM, IgA等）、细胞因子 \\
                \end{tabular}
        \end{table}

        表\ref{tab:compare_between_two_immune_response}总结对比了固有免疫和适应性免疫的不同。另一方面，固有免疫和适应性免疫是有序发生、相互协调的。一方面，固有免疫是适应性免疫的先决条件和启动因素，能提供适应性免疫应答所需的活化信号；另一方面，适应性免疫的效应细胞和分子也可反过来影响和调节固有免疫应答。

\subsection{固有免疫}

    \subsubsection{组织屏障}

    固有免疫系统的组织屏障可抵御病原微生物或异物入侵，主要包括体表皮肤、黏膜屏障和体内屏障。皮肤、黏膜及其附属成分组成的物理、化学和微生物屏障可阻挡和抵御外来病原体或异物的入侵。体内屏障（包括血脑屏障、血胎屏障等）可阻止病原体或异物，以及一些免疫系统成分进入机体重要免疫豁免器官。

        % 表 1.1：皮肤和黏膜屏障
        \begin{table}[H]
            \centering
            \caption{皮肤和黏膜屏障}
            \label{tab:skin_mucosa_barrier}
            \begin{tabular}{@{}lll@{}}
            \toprule
            \textbf{屏障类型} & \textbf{组成/机制} & \textbf{功能/作用} \\
            \midrule
            物理屏障 & 紧密连接的上皮细胞、纤毛定向摆动、黏膜表面分泌液等 & 清除黏膜表面的病原体和异物，阻挡侵入体内 \\
            \midrule
            化学屏障 & 皮肤和黏膜分泌物中的杀/抑菌物质 & 形成抵御病原体感染的化学屏障 \\
            \midrule
            微生物屏障 & 皮肤和黏膜表面的正常菌群 & 与病原体竞争结合上皮细胞、分泌杀/抑菌物质等 \\
            \bottomrule
            \end{tabular}
        \end{table}
        
        % 表 1.2：体内屏障
        \begin{table}[H]
            \centering
            \caption{体内屏障}
            \label{tab:internal_barrier}
            \begin{tabular}{@{}ll@{}}
            \toprule
            \textbf{屏障名称} & \textbf{主要功能} \\
            \midrule
            血脑屏障 & 阻挡血液中的病原体和大分子物质进入脑组织及脑室 \\
            \midrule
            血胎屏障 & 防止母体内病原体及有害物质进入胎儿体内，同时允许营养物质交换 \\
            \midrule
            血-胸腺屏障 & 阻止微生物和大分子物质进入胸腺组织 \\
            \bottomrule
            \end{tabular}
        \end{table}

    \subsubsection{髓系免疫细胞的免疫应答}

    当病原体突破组织屏障后，体内的免疫细胞会被迅速激活并发挥作用。这些细胞通过其表面的模式识别受体（PRRs）来识别病原体上共有的、高度保守的病原相关分子模式（PAMPs）（如细菌的脂多糖、病毒的双链 RNA）。固有免疫细胞主要包括髓系免疫细胞（它们都来源于共同髓系祖细胞）、固有淋巴样细胞（Innate Lymphoid Cells, ILCs）、固有样淋巴细胞（Innate-like Lymphocytes, ILLs）。

    各种髓系免疫细胞同样各有其复杂的应答机制，它们构成一个迅速有效的协作网络。概而论之：
    
        \begin{itemize}
            \item \textbf{早期炎症反应阶段}（分钟至数小时）：肥大细胞作为组织驻留细胞首先被激活，释放组胺等血管活性介质，迅速增加血管通透性，促进血浆蛋白渗出。中性粒细胞作为主要的效应细胞，通过趋化作用迅速迁移至感染部位，执行吞噬作用并清除细菌。
            \item \textbf{炎症级联与免疫调节阶段}（数小时至数天）：单核细胞随后被招募至炎症部位并分化为巨噬细胞，增强吞噬能力并分泌大量促炎细胞因子，放大炎症反应。嗜碱性粒细胞通过分泌 IL-4 和 IL-13 等细胞因子，参与免疫调节，诱导 Th2 型免疫应答。
            \item \textbf{特异性效应与抗原提呈阶段}：在寄生虫感染等特定条件下，嗜酸性粒细胞在 IL-5 的诱导下聚集，释放毒性颗粒蛋白介导细胞毒性作用。同时，树突状细胞摄取并处理抗原，随后迁移至引流淋巴结进行抗原提呈，激活适应性免疫应答。
        \end{itemize}

    \subsubsection{ILCs 和 ILLs 的免疫应答}

    ILCs 来源于骨髓共同淋巴样前体，由转录因子 ID2+ 固有淋巴样前体发育分化而成。包括 ILC1/2/3 三个亚群，以及 NK 细胞。不表达特异性/泛特异性抗原受体，故其活化不依赖于对抗原的识别。此类淋巴细胞表达一系列与其活化或抑制相关的受体，可被感染部位组织细胞产生的某些细胞因子或被某些病毒感染/肿瘤靶细胞表面相关配体激活，并通过分泌不同类的细胞因子参与抗感染免疫、过敏性炎症反应，或通过释放细胞毒性介质使相关靶细胞裂解破坏。
    
    ILLs 来源于骨髓共同淋巴样前体。包括 NKT 细胞 (自然杀伤 T 细胞)，γδT 细胞，B1 细胞。可通过表面有限多样性抗原识别受体直接识别结合病原体感染或肿瘤靶细胞表面某些特定表位分子或某些病原体等抗原性异物而被激活，释放一系列细胞毒性介质使相关靶细胞裂解破坏或产生以 IgM 为主的抗菌抗体发挥抗感染免疫作用。

    \subsubsection{固有免疫分子的应答}

    体液中存在多种可溶性蛋白分子，能直接参与和调控固有免疫应答、清除病原体、介导免疫细胞的信息交流和调控。它们包括补体系统、受体分子、细胞因子等。

    以补体系统为例，补体系统由大约 50 种蛋白质和蛋白片段组成，包括血清蛋白和细胞膜受体，它们约占血清中球蛋白的 10\%。这些蛋白由肝脏合成，并作为无活性的前体在血液中循环。当被几种触发因素刺激时，系统中的蛋白酶会切割特定的蛋白质并启动进一步酶切的扩增级联反应，最终可以刺激吞噬细胞清除入侵微生物，诱发炎症以吸引更多的吞噬细胞参与免疫反应，同时可以形成细胞膜攻击复合物，直接对入侵微生物进行杀伤。

    \subsubsection{炎症反应}

    上述机制共同触发炎症反应过程。炎症反应是机体免疫系统对有害刺激所做出的一种防御性反应，根本目的是清除有害物质、阻止其扩散，并启动组织修复过程。

    \begin{enumerate}
        \item 免疫细胞通过其表面的模式识别受体识别病原体或损伤细胞释放的" 危险信号"
        \item 血管舒张、血管通透性增加
        \item 趋化作用（损伤部位会释放趋化因子，引导血液中的免疫细胞向炎症部位移动）、白细胞渗出（中性粒细胞和单核/巨噬细胞等沿着血管壁滚动、粘附，最终穿过血管壁进入组织）、吞噬作用（到达现场的免疫细胞识别、吞噬并消化病原体或细胞碎片）
        \item 免疫细胞通过吞噬、释放抗菌物质、形成中性粒细胞胞外陷阱等方式消灭病原体；炎症反应后期，抗炎因子开始释放，抑制炎症。成纤维细胞开始产生胶原蛋白，内皮细胞形成新血管，共同促进组织修复和疤痕形成。
    \end{enumerate}

\subsection{适应性免疫应答}

    作为人体针对特定病原体建立的防御机制，适应性免疫高度特异且具有记忆性。如图所示，以 SARS-CoV-2 病毒感染为例，适应性免疫主要分为体液免疫和细胞免疫两条的路径。整个免疫反应始于抗原的识别与呈递。当病毒首次侵入人体后，抗原呈递细胞如巨噬细胞或树突状细胞会吞噬病毒，将其蛋白降解为多肽片段，并结合主要组织相容性复合体 II（MHC II）类分子展示在细胞表面。这些抗原复合物被初始辅助性 T 细胞表面的受体特异性识别，在 CD4 共受体的辅助下，辅助性 T 细胞被激活并增殖，成为调控后续免疫反应的核心枢纽。

    \subsubsection{体液免疫过程}

    \begin{figure}[H]
        \centering
        \includegraphics[width=0.5\linewidth]{figures/humoral-immune-response.jpg}
        \caption{体液免疫过程图示}
        \label{fig:humoral-immunity}
    \end{figure}

    图\ref{fig:humoral-immunity}展示了体液免疫过程，其核心效应细胞是源自骨髓的 B 淋巴细胞。
    \begin{enumerate}
        \item 被激活的辅助性 T 细胞分化为 Th2 细胞及滤泡辅助性 T 细胞，它们通过分泌白细胞介素-4 和白细胞介素-21 等细胞因子，向 B 细胞传递活化信号；
        \item 同时，B 细胞表面的受体直接识别病毒抗原，并在 CD40 配体等共刺激分子的作用下完成活化；
        \item 活化后的 B 细胞发生克隆扩增并分化为浆细胞；
        \item 浆细胞作为效应细胞合成并分泌大量特异性抗体；
        \item 抗体随血液循环到达感染部位，特异性地结合病毒表面的刺突蛋白，从而中和病毒的感染能力，并促进病毒的分解与清除。与此同时，部分 B 细胞分化为生发中心 B 细胞，最终形成长寿的记忆 B 细胞，为再次感染储备免疫记忆；
    \end{enumerate}

    \subsubsection{抗体的结构}

        不同抗体都含有 Ig 单体，称为抗体的基本结构（如\ref{fig:aB-structure}）。Ig 单体是由两条完全相同的相对分子质量较大的重链和两条完全相同的分子质量较小的轻链组成，重链和轻链通过二硫键连接，呈 Y 形结构。每条肽链分别由 2\textasciitilde5 个约含 110 个氨基酸、结构相似但功能不同的结构域（又称功能区）组成。

            \begin{figure}[H]
                \centering
                \includegraphics[width=0.71\columnwidth]{figures/antibody-structure-diagram_683773-264.jpg}
                \caption{抗体结构示意图}
                \label{fig:aB-structure}
            \end{figure}

        Ig单体重链和轻链靠近 N 端约 110 个氨基酸的序列差异很大，其他部分的氨基酸序列则相对恒定。据此，将重链和轻链分为可变区（variable region，V 区）和恒定区（constant region，C 区）。V 区决定了抗体的特异性，是抗体与抗原特异性结合的位点；不同类别抗体的 C 区末端可以与一些自体细胞上的受体结合，介导不同的生物学活性。

        除了这些基本结构外，抗体还有 J 链（joining chain，将多个 Ig 抗体单体连接层抗体复合体）、分泌片（secretory piece，SP，保护抗体免受蛋白酶水解）等其他的辅助成分。

        按照恒定区的种类，人体内的抗体可以分为 IgG，IgA，IgM，IgE，IgD 五种，他们的生物活性、生理功能、存在形式存在着一定的差异。下表\ref{tab:immunoglobulin_compare_compact}比较了这五种不同类型抗体的差异。

        \begin{table}[H]
            \small
            \RaggedRight
            \caption{五类免疫球蛋白（IgG, IgM, IgA, IgD, IgE）比较}
            \label{tab:immunoglobulin_compare_compact}
            \begin{tabular}{l p{2.3cm} p{2.3cm} p{2.3cm} p{2.3cm} p{2.3cm}}
            \toprule
            \textbf{特征} & \textbf{IgG} & \textbf{IgM} & \textbf{IgA} & \textbf{IgD} & \textbf{IgE} \\
            \midrule
            血清含量 & 75\% & 10\% & 15\% & <1\% & <0.01\% \\
            结构形式 & 单体 & 五聚体 & 单体/二聚体 & 单体 & 单体 \\
            辅助成分 & 无 & J链、SP & J链 & 无 & 无 \\
            主要分布 & 血液、组织、胎盘 & 血液、淋巴 & 黏膜、分泌液 & B 细胞表面 & 组织、肥大细胞 \\
            胎盘通过 & 可 & 否 & 否 & 否 & 否 \\
            主要功能 & 中和、调理、ADCC & 初次免疫首发、强补体活化 & 黏膜免疫、母乳抗体 & BCR 成分 & 过敏反应、抗寄生虫 \\
            半衰期 & 最长（21天） & 5天 & 6天 & 3天 & 2天 \\
            分子特性 & 最轻、扩散快 & 最大、五聚体高效结合 & 稳定于黏膜 & - & 结合 Fc$\mathrm{\epsilon}$RI，高敏反应 \\
            \bottomrule
            \end{tabular}
            \tabletext{注：Fc$\mathrm{\epsilon}$RI是肥大细胞、嗜碱性粒细胞上的受体}
        \end{table}

    \subsubsection{抗体的作用方式}


        抗体对病原体有多种作用方式，：V 区主要功能是特异性结合抗原，从而阻断病原入侵，可发挥中和作用；C 区则在 V 区与特异性抗原结合后，通过激活补体及与吞噬细胞表面 Fc 受体结合，发挥抗体依赖的细胞介导的细胞吞噬作用（ADCP）、产生 ADCC 效应、介导超敏反应和通过胎盘等。

        \begin{itemize}
            \item 中和作用（Neutralization）：抗体 V 区与病毒表面的特定结构结合将其封闭，使病原体不能再与宿主细胞结合，无法进入细胞进行繁殖
            \item 调理作用（opsonization）：IgG 抗体头部的 Fab 段与抗原结合，尾部（C 区）Fc 段与巨噬细胞或中性粒细胞表面的 Fc$\gamma$R 结合，诱导吞噬细胞对细菌的吞噬
            \item 抗体依赖的细胞吞噬作用（antibody dependent cellular phagocytosis， ADCC）：抗体与靶细胞表面的抗原特异性结合，随后抗体 Fc 片段与效应细胞（如巨噬细胞）表面 Fc$\gamma$R结合，诱导巨噬细胞吞噬靶细胞，通过吞噬体酸化作用导致靶细胞的内在化和降解
            \item 抗体依赖细胞介导的细胞毒作用（antibody dependent cell mediated cytotoxicity，ADCC）：抗体头部的 Fab 段结合病毒感染的细胞或肿瘤细胞表面的抗原表位，其 Fc 段与杀伤细胞（主要是 NK 细胞）表面的 FcR 结合，介导杀伤细胞直接杀伤靶细胞
            \item 介导I型超敏反应：IgE 为亲细胞抗体，可通过其 Fc 段与肥大细胞和嗜碱性粒细胞表面的高亲和力 IgE Fc 受体（FcεRI）结合，并使其致敏，若之后相同病原体再次进入机体与致敏靶细胞表面特异性 IgE 结合，即可促使这些细胞合成和释放生物活性物质，引起I型超敏反应
        \end{itemize}

        其中调理作用是ADCP的一种宏观陈述，而ADCP更像是最终作用的机理（来自Gemini回复）。I型超敏反应是病理性的反应，最终引发了人体的过敏反应，而其他类型的反应都是保护性的。

        下表\ref{tab:antibody_effector_mechanisms}总结了这些作用方式的不同。

        \begin{table}[H]
            \small
            \RaggedRight
            \caption{抗体五种效应机制的比较}
            \label{tab:antibody_effector_mechanisms}
            \begin{tabular}{l p{2.1cm} p{2.1cm} p{2.1cm} p{2.1cm} p{2.1cm}}
            \toprule
            \textbf{特征} & \textbf{中和作用} & \textbf{调理作用} & \textbf{ADCC} & \textbf{ADCP} & \textbf{I型超敏反应} \\
            \midrule
            作用本质 & 阻断病原体/毒素功能 & 标记病原体，增强识别 & 杀伤靶细胞 & 吞噬并消化靶标 & 过敏介质释放 \\
            \midrule
            关键抗体部位 & Fab段 & Fc段 & Fc段 & Fc段 & Fc段（IgE） \\
            \midrule
            参与抗体类型 & IgG、IgA、IgM & IgG（IgG1、IgG3） & IgG、IgE & IgG & \textbf{仅IgE} \\
            \midrule
            主要效应细胞 & 无需细胞参与 & 巨噬细胞、中性粒细胞 & NK细胞、嗜酸性粒细胞、巨噬细胞 & 巨噬细胞、单核细胞、中性粒细胞 & 肥大细胞、嗜碱性粒细胞 \\
            \midrule
            关键受体 & 无 & Fc$\gamma$R & CD16(Fc$\gamma$RIII)、Fc$\varepsilon$R & Fc$\gamma$RI、Fc$\gamma$RIIA & Fc$\varepsilon$RI \\
            \midrule
            靶标类型 & 游离病毒、毒素、细菌黏附素 & 胞外菌、病毒颗粒 & 感染细胞、肿瘤细胞、寄生虫 & 细菌、小颗粒、凋亡细胞 & 无害的过敏原 \\
            \midrule
            处理方式 & 阻断功能（不清除） & 促进吞噬（协助） & 穿孔素/颗粒酶杀伤 & 吞噬体-溶酶体消化 & 脱颗粒释放介质 \\
            \midrule
            作用速度 & 即时 & 需细胞募集 & 需细胞接触 & 需细胞接触 & \textbf{极快（数分钟）} \\
            \midrule
            典型应用场景 & 病毒感染预防、抗毒素治疗 & 荚膜菌感染（肺炎链球菌） & 病毒感染细胞清除、抗肿瘤 & 细菌清除、免疫复合物清除 & 过敏性鼻炎、哮喘、过敏性休克 \\
            \midrule
            释放物质 & 无 & 无（促进吞噬信号） & 穿孔素、颗粒酶、TNF-$\alpha$ & 溶酶体酶（吞噬后） & 组胺、白三烯、前列腺素、细胞因子 \\
            \bottomrule
            \end{tabular}
            \tabletext{注：Fab即抗原结合片段（fragment of antigen binding），Fc即可结晶片段（fragment crystallizable）}
        \end{table}

        对不同类型的病原体，抗体的类型、作用方式也不同，下表\ref{tab:pathogen_antibody_mechanism}对此进行了简要总结。

        \begin{table}[H]
            \small
            \RaggedRight
            \caption{针对不同病原体的抗体类型与核心作用机制}
            \label{tab:pathogen_antibody_mechanism}
            \begin{tabular}{l p{1.8cm} p{3.2cm} p{4.2cm}}
            \toprule
            \textbf{病原体类型} & \textbf{关键抗体} & \textbf{主要作用机制} & \textbf{免疫效应描述} \\
            \midrule
            \textbf{游离病毒} & IgG, IgA & \textbf{中和作用}、凝集作用 & 阻断病毒与受体结合；IgA在黏膜表面阻止入侵 \\
            \textbf{病毒感染细胞} & IgG & \textbf{ADCC}、补体激活 & 引导NK细胞裂解"病毒工厂"；补体导致细胞溶解 \\
            \textbf{胞外细菌} & IgG, IgM & \textbf{调理作用}、补体激活 & IgG/C3b包被细菌增强吞噬；MAC复合物直接钻孔 \\
            \textbf{细菌外毒素} & IgG & \textbf{中和作用} & 形成抗原-抗体复合物，使毒素失去毒性活性 \\
            \textbf{胞内菌/真菌} & IgG & 调理作用、ADCP & 辅助巨噬细胞吞噬；但在胞内菌中主要依赖细胞免疫 \\
            \textbf{蠕虫/寄生虫} & \textbf{IgE} & \textbf{ADCC} (IgE介导) & 招募嗜酸性粒细胞脱颗粒，释放毒性酶杀伤虫体 \\
            \bottomrule
            \end{tabular}
            \tabletext{注：sIgA=分泌型IgA；MAC=膜攻击复合物；对于胞内感染（如结核、病毒复制期），抗体作用有限，主要依赖T细胞介导的细胞免疫。}
        \end{table}

        \subsubsection{细胞免疫}

    图\ref{fig:cell-meditated-immunity}展示了细胞免疫过程，主要目的是清除已经被病毒感染的宿主细胞。

    \begin{figure}
        \centering
        \includegraphics[width=0.5\linewidth]{figures/cell-immunity.png}
        \caption{细胞免疫过程示意图}
        \label{fig:cell-meditated-immunity}
    \end{figure}

    \begin{enumerate}
        \item 辅助性 T 细胞分化为 Th1 亚群，分泌干扰素-gamma 和白细胞介素-2 等细胞因子，以此激活源自胸腺的细胞毒性 T 细胞；
        \item 被病毒感染的细胞会将病毒抗原通过主要组织相容性复合体 I（MHC I）类分子呈递在表面；
        \item 这一信号被 CD8 阳性的细胞毒性 T 细胞识别，活化后的细胞毒性 T 细胞通过释放穿孔素和颗粒酶，或通过 Fas 配体介导的途径，诱导被感染细胞发生凋亡，从而阻断病毒在细胞内的复制；
        \item 此外，自然杀伤细胞也能在细胞因子的调节下，识别并杀伤病毒感染细胞，起到补充防御的作用。
    \end{enumerate}
    
    免疫记忆的形成是适应性免疫最显著的特征。在初次免疫应答结束后，部分淋巴细胞分化为记忆辅助性 T 细胞、记忆细胞毒性 T 细胞以及记忆 B 细胞，它们处于静息状态长期存活于体内。当机体再次遭遇同种 SARS-CoV-2 抗原时，这些记忆细胞能够迅速被唤醒，并在极短的时间内产生更强烈的体液免疫和细胞免疫反应。这种迅速且高效的二次应答机制，能够在其造成严重病理损伤之前将病原体有效清除，从而赋予机体长期的免疫保护力。

\section{诊断学综述}

    \subsection{查体的通用项目}

    体格检查（physical examination）是指医生运用自己的感官和借助于简便的检查工具，如体温计、血压计、叩诊锤、听诊器、检眼镜等，客观地了解和评估人体状况的一系列最基本的检查方法。许多疾病通过体格检查再结合病史即可作出临床诊断。

    体格检查的基本方法有五种：视诊、触诊、叩诊、听诊和嗅诊。部分专科体格检查还会用到其他一些检查方法，例如骨科检查的动诊和量诊。

    \subsubsection{一般状态的检查}

    一般检查为整个体格检查过程中的第一步，是对病人全身状态的概括性观察，以视诊为主，配合触诊、听诊和嗅诊进行检查。
    
    一般检查除了对患者性别、年龄的判断及基本生命体征的测量外，还有如下的内容：
    \begin{enumerate}
        \item 发育：正常 / 落后 / 营养不良
        \item 体型： 无力型 / 正力型 / 超力型
        \item 营养：良好 / 不良 / 中等
        \item 神志：清楚 / 嗜睡 / 昏睡 / 昏迷
        \item 面容：急性面容、慢性面容、贫血面容等
	    \item 精神：自然 / 烦躁 / 倦怠 / 反应迟钝
        \item 体位：自主体位 / 被动体位 / 强迫体位
        \item 步态：正常 / 步态不稳 / 共济失调步态等
        \item 皮肤黏膜：无黄染 / 无皮疹 / 无出血点 / 无紫癜 / 无脱水体征
        \item 淋巴结：未触及肿大
    \end{enumerate}

    生命体征检查即对患者的体温（T）、脉搏（P）、呼吸（R）、血压（BP）及血氧饱和度 $SpO_2$ 的检查。

    主流的体温检测有腋测法、口测法、肛测法三种，法仅用于体温筛查。体温测定的结果，应按时记录于体温记录单上，描绘出体温曲线。多数发热性疾病，其体温曲线的变化具有一定的规律性。

    对患者脉搏、血压和血氧饱和度的检查常用专门的仪器进行。
    
    正常成人脉率在安静、清醒的情况下为 60\textasciitilde100 次 / 分，老年人偏慢，女性稍快，儿童较快，小于3岁的儿童多在 100 次 / 分以上。各种生理、病理情况或药物影响也可使脉率增快或减慢。此外，除脉率快慢外，还应观察脉率与心率是否一致。某些心律失常如心房颤动或较早出现的期前收缩，由于部分心脏收缩的搏出量低，不足以引起周围动脉搏动，故脉率可少于心率。
    
    正常血压的范围为收缩压 $<120\ mmHg$、舒张压 $<80\ mmHg$，正常高值为收缩压 $120\sim129$ 内、舒张压在 $80\sim89$内，对于更高血压的情况即为高血压，高血压又能分为1\textasciitilde3级。高血压定义为：在未使用降压药物的情况下，非同日 3 次测量诊室血压，$SBP\ge140\ mmHg$ 和 $DBP\ge90\ mmHg$。情绪激动、紧张、运动等因素会导致人体血压测量结果偏高，因此高血压的诊断应结合诊室血压、动态血压监测和家庭自测血压的数据。高血压是动脉粥样硬化、心力衰竭、心房颤动、脑卒中、肾功能不全等的重要危险因素。凡血压低于 90/60mmHg 时称低血压。急性的持续（＞30 分钟）低血压状态多见于严重病症，如休克、心肌梗死、急性心脏压塞等疾病。

    动脉血氧饱和度是判断机体是否缺氧的一个指标，但是反映缺氧并不敏感，而且有掩盖缺氧的潜在危险，一般作为患者肺功能检查的参考指标，参考的正常区间为$95\%\sim98\%$。

    正常成人静息状态下，呼吸为 12\textasciitilde20 次/分，呼吸与脉搏之比为 $1:4$。呼吸过速（tachypnea）指呼吸频率超过 20 次 / 分。见于发热、疼痛、贫血、甲状腺功能亢进及心力衰竭等。一般体温升高 1$^oC$，呼吸频率大约增加 4 次 / 分。 呼吸过缓（bradypnea） 指呼吸频率低于 12 次/分。呼吸浅慢见于麻醉剂或镇静剂过量和颅内压增高等。

    淋巴结分布于全身，一般体格检查即检查身体各部表浅的淋巴结是否肿大。检查淋巴结的方法是视诊和触诊。视诊时不仅要注意局部征象（包括皮肤是否隆起，颜色有无变化，有无皮疹、瘢痕、瘘管、窦道等）也要注意全身状态。

    体表淋巴结触及可能有未触及、局部性淋巴结肿大、全身性淋巴结肿大三种情况。局部性淋巴结肿大常指示非特异性淋巴结炎（由引流区域的急、慢性炎症所引起）、单纯淋巴结炎（淋巴结本身的急性炎症）、淋巴结核（肿大的淋巴结相互粘连或与周围组织粘连）、恶性肿瘤淋巴结转移等病症，全身性淋巴结肿大常指示一些感染性疾病（如传染性单核细胞增多症、艾滋病等病毒感染和结核、布鲁氏病菌等细菌感染及其他的螺旋体、原虫与寄生虫感染等）和结缔组织疾病（如系统性红斑狼疮、干燥综合征、结节病等）、血液系统疾病（如急、慢性白血病，淋巴瘤，恶性组织细胞病等）等非感染性疾病。

    面容是指面部呈现的状态，表情是在面部或姿态上思想感情的表现。健康人表情自然，神态安怡。患病后因病痛困扰，常出现痛苦、忧虑或疲惫的面容与表情。某些疾病发展到一定程度时，尚可出现特征性的面容与表情，对疾病的诊断具有重要价值。

    除急性面容外，还有很多其他的病态面容。慢性病容主要表现为面容憔悴，面色晦暗或苍白无华，目光暗淡、表情忧虑。见于慢性消耗性疾病，如恶性肿瘤、肝硬化、严重结核病等。贫血面容表现为面色苍白，唇舌色淡，表情疲惫。见于各种原因所致的贫血。肾疾病、肝疾病、甲状腺亢进及减退患者的面容也有对应的表现，还有伤寒面容、面具面容等多种面容，指示不同的疾病，为疾病进一步检查和诊断提供参考。

    
    \subsubsection{头部检查}

    本案例中进行的头部检查是咽喉部的检查。

    咽部的检查方法是，病人取坐位，头略后仰，口张大并发"啊"音，此时医生用压舌板在舌的前 2/3 与后 1/3 交界处迅速下压，此时软腭上抬，在照明的配合下即可见软腭、腭垂、软腭弓、扁桃体、咽后壁等。

    检查时若发现咽部黏膜充血、红肿、黏膜腺分泌增多，多见于急性咽炎；若咽部黏膜充血、表面粗糙，并可见淋巴滤泡呈簇状增殖，见于慢性咽炎。扁桃体发炎时，腺体红肿、增大，在扁桃体隐窝内有黄白色分泌物，或渗出物形成的苔片状假膜，很易剥离，这点与咽白喉在扁桃体上所形成的假膜不同，白喉假膜不易剥离，若强行剥离则易引起出血。
    
    如图\ref{fig:biantaotifenji}所示，扁桃体增大一般分为三度：不超过腭咽弓者为$I$度；超过腭咽弓者为$II$度；达到或超过咽后壁中线者为$III$度。一般检查未见扁桃体增大时可用压舌板刺激咽部，引起反射性恶心，如看到扁桃体突出为包埋式扁桃体，同时隐窝有脓栓时常构成反复发热的隐性病灶。

    \begin{figure}[H]
        \centering
        \includegraphics[width=0.5\linewidth]{figures/bttfenji.png}
        \caption{扁桃体三度分级示意图}
        \label{fig:biantaotifenji}
    \end{figure}

    \subsubsection{胸部检查}

    本案例中进行的胸部检查有呼吸音和心音的检查。

    呼吸音听诊时受检者取坐位或卧位，听诊的顺序一般由肺尖开始，自上而下分别检查前胸部、侧胸部和背部。听诊前胸部时沿锁骨中线和腋前线；听诊侧胸部时沿腋中线和腋后线；听诊背部时沿肩胛线，自上至下逐一肋间进行，而且要在上下、左右对称的部位进行对比。

    病例中出现的粗糙性呼吸音为支气管黏膜轻度水肿或炎症浸润造成不光滑或狭窄，使气流进出不畅所形成的粗糙呼吸音，常见于支气管或肺部炎症的早期。

    啰音（crackles，rales）是呼吸音以外的附加音，该音正常情况下并不存在。啰音又分为干啰音和湿啰音两种。（1）湿啰音（moist crackles） 系由于吸气时气体通过呼吸道内的分泌物如渗出液、痰液、血液、黏液和脓液等，形成的水泡破裂所产生的声音，肺部局限性湿啰音，仅提示该处的局部病变，如肺炎、肺结核或支气管扩张等。两侧肺底湿啰音，多见于心力衰竭所致的肺淤血和支气管肺炎等。如两肺野满布湿啰音，则多见于急性肺水肿和严重支气管肺炎。（2）干啰音（wheezes，rhonchi） 系由于气管、支气管或细支气管狭窄或部分阻塞，空气吸入或呼出时形成湍流所产生的声音。呼吸道狭窄或不完全阻塞的病理基础包括炎症引起的黏膜充血水肿和分泌物增加；支气管平滑肌痉挛；管腔内肿瘤或异物阻塞；以及管壁被管外肿大的淋巴结或纵隔肿瘤压迫引起的管腔狭窄等。发生于双侧肺部的干啰音，常见于支气管哮喘、慢性支气管炎、慢性阻塞性肺疾病和心源性哮喘等。局限性干啰音，是局部支气管狭窄所致，常见于支气管内膜结核或肿瘤等。
    
    病例中患者常年吸烟、身患慢性阻塞性肺炎且未规律治疗，查体时未闻及干湿啰音可能是由于慢阻肺疾病并不严重。符合感冒引起的呼吸道炎症早期的特征。

    心音（heart sound）是心动周期内高压血液冲击心脏内的瓣膜使其振动而发出的声音，按其在心动周期中出现的先后次序，可依次命名为第一心音（first heart sound，S1）、第二心音（second heart sound，S2）、第三心音（third heart sound，S3）和第四心音（fourth heart sound，S4）。正常人一般只能听到第一、第二心音。第三心音多为病理性心音，但也可在部分青少年中闻及。第四心音属病理性心音。

    除正常节律性心音外，心脏听诊还可以听到额外心音（extra cardiac sound） 指在正常 S1、S2 之外听到的附加心音，与心脏杂音不同。多数为病理性，大部分出现在 S2 之后即舒张期，与原有的心音 S1、S2 构成三音律（triple rhythm），如奔马律、开瓣音和心包叩击音等；也可出现在 S1 之后即收缩期，如收缩期喷射音。

    心脏杂音（cardiac murmurs）是指除心音与额外心音外，在心脏收缩或舒张期发现的异常声音，杂音性质的判断对于心脏病的诊断具有重要的参考价值。

    在本病例中，各瓣膜听诊区未闻及杂音。

    \subsubsection{腹部检查}

    本病例中进行的腹部检查为腹部的视诊和腹壁紧张度、肝脾肋下触诊和肝肾的扣诊。

    叩痛（tenderness on percussion）是查体中常用于判断实质脏器是否存在炎症、肿大、管腔梗阻、包膜张力升高或被膜牵拉的体征。其意义在于辅助快速筛查内脏炎症（如肝炎、胆囊炎、肾盂肾炎）、辅助诊断内脏肿大（如肝肿大、肾肿大）、评估病变侧性（左或右肾区痛判断病灶）并指导进一步检查（如腹部超声、CT、血、尿检查）。肝区叩痛可能指示各类肝炎、肝脓肿、肝充血性肿大等病症，肾区叩痛常指示急性肾盂肾炎、肾结石或输尿管结石、肾周围感染或肾周脓肿、多囊肾、肾肿瘤等导致肾脏增大等疾病。在本病例中，患者肝肾叩痛反应呈阴性。

    正常人腹壁有一定张力，但触之柔软，较易压陷，称腹壁柔软，有些人（尤其儿童）因不习惯触摸或怕痒而发笑致腹肌自主性痉挛，称肌卫增强，在适当诱导或转移注意力后可消失，不属异常。某些病理情况可使全腹或局部腹肌紧张度增加或减弱。在本病例中，患者触诊腹平软，属于正常状况。

    正常腹部触摸时不引起疼痛，重按时仅有一种压迫感。真正的压痛（tenderness）多来自腹壁或腹腔内的病变。腹腔内的病变，如脏器的炎症、淤血、肿瘤、破裂、扭转以及腹膜的刺激（炎症、出血等）等均可引起压痛，压痛的部位常提示存在相关脏器的病变。当医生用手触诊腹部出现压痛后，用并拢的 2\textasciitilde3 个手指（示、中、环指）压于原处稍停片刻，使压痛感觉趋于稳定，然后迅速将手抬起，如此时病人感觉腹痛骤然加重，并常伴有痛苦表情或呻吟，称为反跳痛（rebound tenderness）。反跳痛是腹膜壁层已受炎症累及的征象，是当突然抬手时腹膜被激惹所致，是腹内脏器病变累及邻近腹膜的标志。在本病例中，患者均未出现压痛及反跳痛。

    腹腔内重要脏器较多，病理情况下常可触到脏器增大或局限性肿块，对诊断有重要意义。肝脏触诊主要用于了解肝脏下缘的位置和肝脏的质地、表面、边缘及搏动等。正常情况下脾脏不能触及。在本病例中，肝脾于肋下未触及，说明患者肝脾大小正常、未出现肿大。

    \subsection{血常规检查}

    血液常规检测（blood routine test）一般指包括红细胞、白细胞、血小板计数、血红蛋白浓度测定及形态学观察等的检查，随着检验技术的发展，快速、自动化、多指标联合的血液分析仪器已广泛应用，血液常规项目的内涵也在逐渐增多，除有形成分计数外，红细胞个体形态、血红蛋白状态、网织红细胞定量及分级、血小板个体形态、白细胞自动分类及异常白细胞提示，甚至外周血有核红细胞数量都已逐渐成为常规检测内容，因此也有把血液常规检测称为全血细胞计数（complete blood count，CBC）。

    红细胞计数（red blood cell count）是指每升全血中红细胞的个数，在正常情况下常于血红蛋白浓度（hemoglobin，Hb 或 HGB，每升全血中血红蛋白含量）结合对机体生成红细胞能力进行判断并能协助诊断与红细胞有关的疾病。红细胞及血红蛋白数量升高通常是因为血浆容量减少（腹泻、呕吐等）或红细胞绝对数量增多（指示缺氧等因素引起的代偿性增多、肿瘤或肾脏疾病或真性红细胞增多症等），红细胞和血红蛋白数量下降常指示各种类型的贫血。在本病例中，红细胞检测值为正常值。

    白细胞计数包括对中性粒细胞、嗜酸性粒细胞、嗜碱性粒细胞、淋巴细胞和单核细胞的计数。（1）中性粒细胞数量增多常指示急性感染（尤其是细菌感染）、严重的组织损伤及大量血细胞破坏、急性出血、急性中毒、各种肿瘤的发生等病症，数量减少指示革兰氏阴性杆菌感染（如伤寒、副伤寒杆菌感）、血液系统疾病、自身免疫疾病等病症。中性粒细胞核分叶数目的左移常指示细菌性感染，右移指示巨幼细胞贫血及造血功能衰退，也可能由抗代谢药物引起（2）嗜酸性粒细胞数量增多常指示过敏性疾病、寄生虫病、血液皮肤疾病和肿瘤等，减少常见于伤寒、副伤寒初期，大手术、烧伤等应激状态，临床意义不大。（3）嗜碱性粒细胞增多指示过敏性疾病、血液病和恶性肿瘤等病症，减少无临床意义。（4）病理性淋巴细胞数量增多常见于感染性疾病、成熟淋巴细胞肿瘤、移植排斥反应等病症，数量减少主要见于应用肾上腺皮质激素、烷化剂、抗淋巴细胞球蛋白等的治疗以及放射线损伤、T 淋巴细胞免疫缺陷病、丙种球蛋白缺乏症（B 淋巴细胞免疫缺陷）等。（5）单核细胞数量增多常指示一些感染和血液疾病，数量减少一半没有临床意义，长期数量的持续减少可能指示病毒感染、血液系统疾病、风湿性关节炎、系统性红斑狼疮等，也可见于一些肿瘤、艾滋病等等疾病。

    血小板数量增多常见于原发性增多见于骨髓增殖性肿瘤，如真性红细胞增多症、原发性血小板增多症、原发性骨髓纤维化早期及慢性髓系白血病等；反应性增多见于急性感染、急性溶血、某些恶性肿瘤。其数量减少可见于血小板生成障碍（如再生障碍性贫血、放射性损伤、急性白血病、巨幼细胞贫血、骨髓纤维化晚期）、血小板破坏或消耗增多（如原发免疫性血小板减少症（ITP）、系统性红斑狼疮（SLE）、淋巴瘤、上呼吸道感染等）、血小板分布异常（如脾肿大、肝硬化等）。

    在本病例中，患者血常规结果中白细胞总数正常说明没有明显外周白细胞反应性增高，这与怀疑的病毒性感染或轻度感染早期症状相符，中性粒正常提示暂时未见明显继发细菌感染，血小板、血红蛋白正常证明无明显出血、贫血或骨髓抑制，淋巴细胞减少可能源于淋巴细胞向感染部位或淋巴器官重分布或滞留、病毒诱导的淋巴细胞凋亡/耗竭等，提示患者目前可能经历病毒性感染引起的"相对淋巴细胞减少"。

    \subsection{感染免疫学实验检查}

    本病例还进行了C-反应蛋白（C-reactive protein）的检查。

    C 反应蛋白（C-reactive protein，CRP）是一种由肝脏合成的，能与肺炎双球菌细胞壁 C 多糖起反应的急性时相反应蛋白，广泛存在于血清和其他体液中。CRP 不仅能结合多种细菌、真菌及原虫等体内的多糖物质，而且在钙离子存在下，还可以结合卵磷脂和核酸等，具有激活补体、促进吞噬和调节免疫的作用。

    CRP 是急性时相反应极灵敏的指标。CRP 升高常见于化脓性感染、组织坏死（心肌梗死、严重创伤、大手术、烧伤等）、恶性肿瘤、结缔组织病、器官移植急性排斥等。还能鉴别细菌性或非细菌性感染，前者 CRP 升高，后者不升高。能够鉴别风湿热活动期和稳定期，前者升高，后者不升高。鉴别器质性和功能性疾病：前者升高，后者不升高。


\section{甲型流感病毒}

\subsection{形态与结构}

    流感病毒在细胞培养中多呈球形，直径为 80-120nm。新分离的流感病毒呈多形态性，以丝状多见。流感病毒的结构由核衣壳和包膜组成。
    
    \begin{enumerate}
        \item 核衣壳
        \begin{itemize}
            \item 分节段的单负链 RNA (-ssRNA)
            \item 与-ssRNA 结合的核蛋白NP
            \item RNA 依赖的 RNA 聚合酶RdRp，包括 PA、PB1、PB2 三个亚基
        \end{itemize}
        甲型和乙型流感病毒的基因组分为 8 个节段（segment），丙型流感病毒基因组分为 7 个节段。基因组的多数节段只编码一种蛋白质，每个节段的两端 12\textasciitilde13 个核苷酸为高度保守序列，与病毒复制关系密切。
        \item 包膜
        \begin{itemize}
            \item 外层：主要来自宿主细胞的脂质双层膜，表面分布呈放射状排列的两种刺突，即血凝素（hemagglutinin，HA）和神经氨酸酶（neuraminidase，NA）；HA 和 NA 均为糖蛋白，比例约为 5ff1。此外，包膜上还分布有基质蛋白（M）
            \begin{enumerate}
                \item HA：由三条糖蛋白链组成的三聚体，HA 的原始肽链 HA0 在蛋白酶作用下裂解肽链中的精氨酸，形成二硫键连接的 HA1 和 HA2 后才具有感染性；这一蛋白酶只存在于呼吸道，决定了流感病毒感染的组织特异性
                \item NA：由四条糖蛋白链组成的四聚体，具有神经氨酸酶活性参与病毒释放、促进病毒扩散，能诱导机体产生特异性抗体
                \item M2 蛋白：贯穿病毒包膜，具有离子通道作用，可使包膜内 pH 下降，有助于病毒进入细胞
            \end{enumerate}
            \item 内层：基质蛋白层，主要由 M1 蛋白构成，能保护病毒核心、维持病毒形态。
        \end{itemize}
    \end{enumerate}

\subsection{基因组}

        \begin{table}[H]
            \centering
            \caption{甲型流感病毒基因组结构}
            \label{tab:influenza_genome}
            \begin{tabular}{@{}llll@{}}
            \toprule
            \textbf{基因节段} & \textbf{核苷酸数} & \textbf{编码的蛋白质} & \textbf{蛋白质功能} \\
            \midrule
            1 & 2 341 & PB2 & -- \\
            2 & 2 341 & PB1 & RNA 聚合酶组分 \\
            3 & 2 233 & PA & -- \\
            4 & 1 778 & HA & 血凝素，为包膜糖蛋白，介导病毒吸附，酸性情况下介导膜融合 \\
            5 & 1 565 & NP & 核蛋白，为病毒衣壳成分，参与病毒转录和复制 \\
            6 & 1 413 & NA & 神经氨酸酶，促进病毒释放 \\
            7 & 1 027 & M1 & 基质蛋白，促进病毒装配 \\
            8 & 890 & M2 & 膜蛋白，为离子通道，促进脱壳 \\
            9 & 890 & NS1 & 非结构蛋白，抑制 mRNA 前体的拼接，拮抗干扰素作用 \\
            10 & 890 & NS2 & 非结构蛋白，帮助病毒 RNP 出核 \\
            \bottomrule
            \end{tabular}
        \end{table}

\subsection{血清型、亚型及变异}

\paragraph{血清型} 根据甲（A）、乙（B）、丙（C）型流感病毒属各自的 NP 和 M1 蛋白抗原性差异，可将流感病毒分为甲（A）、乙（B）、丙（C）血清型（serotype）。

\paragraph{亚型} 甲型流感病毒的 HA 和 NA 均容易发生变异，根据 HA 和 NA 的抗原性差异，可将甲型流感病毒分为若干亚型（subtype），目前 HA 包括 16 种亚型（H1-H16），NA 包括 9 种亚型（N1-N9）。乙型流感病毒抗原存在一定变异，但尚无亚型之分；丙型流感病毒至今也无亚型。

\paragraph{变异}流感病毒特别是甲型流感病毒容易发生抗原变异。甲型流感病毒变异除病毒 RdRp 缺乏校对（proof-reading）机制，导致基因组复制中易形成突变并被保留下来外，还因为分节段RNA 基因组容易发生基因重排（gene reassortment）。根据甲型流感病毒抗原性变异的程度，分为两种形式：
\begin{enumerate}
    \item 抗原漂移（antigenic drift）：抗原变异幅度小，主要是 HA 氨基酸的变异，其次是 NA 氨基酸的变异，属于量变。这种变异是由病毒基因点突变导致氨基酸变异，并可不断累积，由于人群存在一定的免疫保护作用，不会引起流感大流行，仅引起中、小规模流行，多出现在寒冷的季节，称为季节性流感。
    \item 抗原转换（antigenic shift）：指抗原变异幅度大，HA 或 NA 氨基酸的变异率达到 30\%\textasciitilde50\%，属于质变，常形成新的亚型（如 H1N1→H2N2）。这种变异多由基因重排形成的新亚型。如人群对流感病毒新亚型易感，并普遍缺乏针对新亚型的免疫力，可导致流感大流行。
\end{enumerate}

\subsection{复制}


    人类流感病毒通过呼吸道途径传播，病毒进入人体呼吸道后，会靶向呼吸道或肠道上皮细胞进行感染和有效复制。此外，部分禽流感甲型病毒（尤其是H7亚型）已被证实可导致人类眼部感染和结膜炎（结膜组织炎症）。人类感染严重程度与病毒在下呼吸道的复制相关，该过程伴随免疫细胞浸润引发的严重炎症反应。

    在细胞层面，流感病毒常在呼吸道上皮细胞内复制，流感病毒首先依附结合到靶细胞的表面，这种结合由病毒血凝素（HA）介导（该蛋白能与细胞表面糖蛋白寡糖中的唾液酸结合），同时也是病毒颗粒与红细胞共孵育时引发血凝反应的原因。人类流感病毒的HA优先识别通过$\alpha2,6$糖苷键连接寡糖链的唾液酸，结合完成后，病毒以内吞体形式被内化，随后内吞体在细胞内转运并酸化，触发病毒HA构象改变，从而诱导病毒包膜与内吞体膜融合。由于不同宿主体内内吞体pH值存在差异，HA的pH稳定性成为病毒趋向性的决定因素之一。病毒包膜与内吞体膜融合后，病毒会释放其内部的八个RNA片段释放到靶细胞细胞质内。随后，病毒核糖核蛋白复合体（vRNPs）被导入受感染细胞的细胞核中。
    
    在核内，借助附着于vRNPs的病毒聚合酶复合物的酶活性，病毒RNA开始进行转录和复制。病毒RNA复制通过一个正链中间体——互补核糖核蛋白复合体（cRNP）完成。病毒RNA转录产生带帽并聚腺苷酸化的正链mRNA，这些mRNA被输出至细胞质，翻译成病毒蛋白。新合成的病毒聚合酶（PB1、PB2和PA）与病毒核蛋白（NP）被运回细胞核，以加速病毒RNA的合成；而病毒膜蛋白HA、NA和M2则转运并嵌入质膜。新合成的HA蛋白需经细胞蛋白酶切割成HA1和HA2多肽方能具备功能，HA的裂解位点同时决定了病毒的组织趋向性。除高致病性禽流感病毒外，所有流感病毒的HA裂解位点均能被呼吸道和肠道上皮细胞中的细胞外蛋白酶识别。
    
    病毒非结构蛋白NS1、PB1-F2（可形成二聚体或多聚体）及PA-x通过调控细胞进程来瓦解宿主抗病毒反应。在病毒感染后期，病毒基质蛋白M1和核输出蛋白NEP定位于细胞核，与病毒核糖核蛋白复合体结合并介导其向细胞质输出。在细胞质中，这些复合体通过循环内体相互作用迁移至质膜，最终组装成八个病毒核糖核蛋白复合体束。新生病毒粒子的出芽过程随之发生，将病毒核糖核蛋白复合体整合至新的病毒颗粒中——该颗粒的膜结构源自宿主细胞质膜，并嵌有病毒跨膜蛋白。神经氨酸酶活性可阻止新生病毒粒子的血凝素与病毒糖蛋白及感染细胞膜上含唾液酸的受体发生无效结合，从而促进病毒传播。病毒复制会导致细胞死亡并引发病理学后果。另外，病毒侵染的靶细胞复合体产生的物质会诱导促炎反应，引发招募先天性和适应性免疫细胞以清除病毒，但过度反应则会引发免疫病理学和肺炎。

            \begin{figure}[H]
                \centering
                \includegraphics[width=0.71\columnwidth]{figures/InfluenzaLifeSpan.png}
                \caption{流感病毒的生命周期}
                \label{fig:figure}
            \end{figure}

\subsection{培养}

    分离培养流感病毒最常用的方法是鸡胚培养。初次分离以羊膜腔接种为宜，传代培养则采用尿囊腔接种。流感病毒在鸡胚增殖后不引起明显的病理改变，常需用血凝试验检测流感病毒并判定其效价。细胞培养常用狗肾传代细胞和猴肾细胞，但流感病毒增殖后引起的CPE不明显，需用红细胞吸附试验（hemoadsorption test）或免疫荧光方法来判定流感病毒的感染和增殖情况。

\subsection{对免疫系统的影响}    

    侵入人体后，流感病毒会攻击人体的免疫系统，使之功能紊乱或异常。其核心原理是流感病毒能够系统性地抑制和规避宿主的先天免疫反应，特别是干扰素 (IFN) 系统 。

    正常情况下，当病毒入侵时，人体的免疫细胞会通过 RIG-I 等"哨兵"蛋白（PAMPs识别受体）来识别病毒的RNA 。这种识别会启动一系列信号传导（通过MAVS蛋白等），最终激活 IRF3、IRF7 和 NF-κB 等关键转录因子，促使细胞产生干扰素 。NS1是流感病毒最主要的免疫抑制蛋白，它会靶向宿主的TRIM25和Riplet (RNF135)等蛋白，阻止它们激活RIG-I，甚至还能直接与 RIG-I 结合，抑制其功能 。PB一族的蛋白能够在MAVS 层面抑制干扰素的产生。即便细胞成功启动了反应，流感病毒也会阻止抗病毒信号进入免疫细胞细胞核引发进一步的反应。
    
    一般情况下，在机体成功产生并释放干扰素 (IFN) 后，IFN 会作为传递病毒信息给相关区域邻近的细胞 。邻近细胞通过其表面的干扰素受体(IFNAR)接收信号，并激活 JAK/STAT 信号通路 ，从而表达数百种抗病毒基因 (ISGs) ，建立起机体"抗病毒"的状态。流感病毒侵染细胞释放的HA、PB2、NS1等蛋白能够作用于这一通路，进而避免引发更大规模的免疫反应。


\section{流行性感冒}

\subsection{流行病学}

    \begin{enumerate}
    \item 传染源：患者和隐性感染者为主要传染源。从潜伏期末到急性期都有传染性, 病毒在人呼吸道分泌物中一般持续排毒 3-7 天, 儿童、免疫功能受损及危重患者病毒排毒时间可超过 1 周。
    \item 传播途径：流感病毒主要通过飞沫传播, 患者呼吸道排出的飞沫播散到周围空气, 被健康人群吸入感染。也可以经气溶胶传播。还可以通过直接或间接接触被病毒污染的手、物品等接触传播。
    \item 人群易感性：人群普遍易感, 感染后对同一亚型会获得一定程度的免疫功能, 但不同亚型间无交叉免疫, 故人体可反复患病。
    \item 流行特征：流感病毒传入人群后, 具有较强传染性, 且抗原极易发生变异, 加之以呼吸道传播为主, 极易引起流行和大流行。流行往往突然发生, 迅速蔓延, 发病率高和流行过程持续时间短是流感的流行特征。流感大流行无明显季节性。在我国温带或寒温带地区, 流感的散发流行一般多发生于冬春季, 在亚热带地区或热带地区, 则更多是在夏季。一般流感流行在我国北方重于南方。流感在人群中蔓延的速度和广度与人口密度有关。流行后人群重新获得一定的免疫功能。
    \end{enumerate}

\subsection{常见病症及其原理}

    流行性感冒的症状主要以发热、头痛、肌肉关节酸痛为主，发热时体温可达39\textasciitilde40$^o\mathrm{C}$，常有咽喉痛、干咳、鼻塞、流涕等，可有畏寒、寒战、乏力、食欲减退等全身症状，部分患者症状轻微或无症状。流感病毒感染可导致慢性基础疾病加重。儿童的发热程度通常高于成人，乙型流感患儿恶心、呕吐、腹泻等消化道症状也较成人多见。新生儿可仅表现为嗜睡、拒奶、呼吸暂停等。

    老年人的临床表现可能不典型，常无发热或为低热，咳嗽、咳痰、气喘和胸痛明显。也可表现为厌食和精神状态改变。无并发症者病程呈自限性，多于发病 3\textasciitilde5 天后发热逐渐消退，全身症状好转，但咳嗽、体力恢复常需较长时间。

    \paragraph{发热} 发热（fever）是指机体在致热原（pyrogen）作用下或各种原因引起体温调节中枢的功能障碍时，体温升高超出正常范围。正常人的体温受体温调节中枢调控，并通过神经、体液因素使产热和散热过程呈动态平衡，保持体温在相对恒定的范围内。正常人体温一般为 36\textasciitilde37℃，可因测量方法不同而略有差异。正常体温在不同个体之间略有差异，且常受机体内、外因素的影响稍有波动。在 24 小时内下午体温较早晨稍高，剧烈运动、劳动或进餐后体温也可略升高，但一般波动范围不超过 1℃。老年人因代谢率偏低，体温相对低于青壮年。另外，在高温环境下体温也可稍升高。

    在正常情况下，人体的产热和散热保持动态平衡。由于各种原因导致产热增加或散热减少，则出现发热。致热原包括外源性和内源性两大类。

    1. \textbf{外源性致热原}（exogenous pyrogen） 外源性致热原的种类甚多，包括：
    \begin{enumerate}
            \item 各种微生物病原体及其产物，如细菌、病毒、真菌及细菌毒素等；
            \item 炎性渗出物及无菌性坏死组织；
            \item 抗原抗体复合物；
            \item 某些类固醇物质，特别是肾上腺皮质激素的代谢产物原胆烷醇酮（etiocholanolone）；
            \item 多糖体成分及多核苷酸、淋巴细胞激活因子等。
    \end{enumerate}

    外源性致热原多为大分子物质，特别是细菌内毒素分子量非常大，不能通过血脑屏障直接作用于体温调节中枢，而是通过激活血液中的中性粒细胞、嗜酸性粒细胞和单核巨噬细胞系统，使其产生并释放内源性致热原，通过下述机制引起发热。

    2. \textbf{内源性致热原}（endogenous pyrogen）又称白细胞致热原（leukocytic pyrogen），如白介素（IL-1）、肿瘤坏死因子（TNF）和干扰素等。内源性致热原一方面通过血-脑脊液屏障直接作用于体温调节中枢的体温调定点（setpoint），使调定点（温阈）上升，体温调节中枢必须对体温加以重新调节发出冲动，并通过垂体内分泌因素使代谢增加或通过运动神经使骨骼肌阵缩（临床表现为寒战），使产热增多；另一方面可通过交感神经使皮肤血管及竖毛肌收缩，停止排汗，散热减少。这一综合调节作用使产热大于散热，体温升高引起发热。非致热原性发热 常见于以下几种情况：
    \begin{enumerate}
            \item 体温调节中枢直接受损：如颅脑外伤、出血、炎症等。
            \item 引起产热过多的疾病：如癫痫持续状态、甲状腺功能亢进症等。
            \item 引起散热减少的疾病：如广泛性皮肤病变、心力衰竭等。
    \end{enumerate}

    发热的病因很多，临床上可分为感染性与非感染性两大类，而以前者多见。

    感染性发热（infective fever）各种病原体如病毒、细菌、支原体、立克次体、螺旋体、真菌、寄生虫等引起的感染，不论是急性、亚急性或慢性，局部性或全身性，均可出现发热。非感染性发热（noninfective fever）主要有下列几类病因：
    
    \begin{enumerate}
    \item 血液病：如白血病、淋巴瘤、恶性组织细胞病等。
    \item 结缔组织疾病：如系统性红斑狼疮、皮肌炎、硬皮病、类风湿关节炎和结节性多动脉炎等。
    \item 变态反应性疾病：如风湿热、药物热、血清病、溶血反应等。
    \item 内分泌代谢疾病：如甲状腺功能亢进症、甲状腺炎、痛风和重度脱水等。
    \item 血栓及栓塞疾病：如心肌梗死、肺梗死、脾梗死和肢体坏死等，通常称为吸收热。
    \item 颅内疾病：如脑出血、脑震荡、脑挫伤等，为中枢性发热。癫痫持续状态可引起发热，为产热过多所致。7. 皮肤病变：皮肤广泛病变致皮肤散热减少而发热，见于广泛性皮炎、鱼鳞病等。慢性心力衰竭使皮肤散热减少也可引起发热。
    \item 恶性肿瘤：各种恶性肿瘤均有可能出现发热。
    \item 物理及化学性损害：如中暑、大手术后、内出血、骨折、大面积烧伤及重度安眠药中毒等。
    \item 自主神经功能紊乱：由于自主神经功能紊乱，影响正常的体温调节过程，使产热大于散热，体温升高，多为低热，常伴有自主神经功能紊乱的其他表现，属功能性发热范畴。常见的功能性低热有：
        \begin{enumerate}
        \item 原发性低热：由于自主神经功能紊乱所致的体温调节障碍或体质异常，低热可持续数月甚至数年之久，热型较规则，体温波动范围较小，多在 0.5℃ 以内。
        \item 感染治愈后低热：由于病毒、细菌、原虫等感染致发热后，低热不退，而原有感染已治愈。此系体温调节功能仍未恢复正常所致，但必须与因机体抵抗力降低导致潜在的病灶（如结核）活动或其他新感染所致的发热相鉴别。
        \item 夏季低热：低热仅发生于夏季，秋凉后自行退热，每年如此反复出现，连续数年后多可自愈。多见于幼儿，因体温调节中枢功能不完善，夏季身体虚弱，且多发生于营养不良或脑发育不全者。
        \item 生理性低热：如精神紧张、剧烈运动后均可出现低热。月经前及妊娠初期也可有低热现象。
        \end{enumerate}
    \end{enumerate}

    \paragraph{呼吸道症状} 咳嗽（cough）与咳痰（expectoration）是临床最常见的呼吸道症状之一。咳嗽是一种反射性防御动作，通过咳嗽可以清除呼吸道内分泌物或异物。但是咳嗽也有不利的一面，如咳嗽可使呼吸道感染扩散，剧烈咳嗽可诱发咯血及自发性气胸等。因此，如果频繁咳嗽影响工作与休息，则为病理状态。痰液是气管、支气管、肺泡内的分泌物，借助咳嗽将其排出称为咳痰。

    \paragraph{其他症状} 流涕也是感冒常见的症状，其原因是鼻黏膜炎症与分泌亢进导致的。感冒也常伴随着咽痛，通常是炎症反应刺激咽部痛觉神经导致。肌肉酸痛、头痛等全身症状也与炎症反应引起的细胞因子风暴有关。

\subsection{诊断}

    流感的诊断需综合流行病学史、临床表现和病原学检查进行判断。在流感流行季节，即使临床表现不典型，尤其是存在流感重症高危因素或需住院治疗的患者，应首先考虑流感的可能性，并及时进行病原学检测。在流感散发季节，对于疑似病毒性肺炎的住院患者，除检测常见呼吸道病原体外，亦应同步开展流感病毒检测。
    
    诊断时应重点关注具有明确流行病学史的患者，包括发病前7天内在无有效个人防护情况下与疑似或确诊流感患者有密切接触者、属于流感样病例聚集发病者之一，或有明确传染他人证据者。同时需结合患者是否具备典型流感临床表现进行综合判断。
    
    实验室检查方面，病毒性呼吸道感染患者的血常规结果通常表现为白细胞计数正常或降低，中性粒细胞比例下降，淋巴细胞比例或绝对值升高，C-反应蛋白正常或仅轻度升高。值得注意的是，部分甲型流感患者可能出现淋巴细胞计数下降，这可能与流感病毒对机体免疫系统的抑制作用有关。上述表现需与细菌性呼吸道感染相鉴别，后者通常表现为白细胞及中性粒细胞计数显著升高，C-反应蛋白明显增高。临床实践中还需警惕病毒与细菌混合感染的可能性。此外，常规检查还应涵盖支原体等非典型病原体的筛查，一般通过相应的IgM抗体检测进行判断。
    
    病原学检查是确诊流感的金标准。核酸检测具有较高的敏感性和特异性，临床应根据流行情况合理选择检测项目，通常需同时进行甲型流感病毒、乙型流感病毒及新型冠状病毒的核酸检测，以明确病原体类型并指导后续治疗。
    
    对于病情较重的患者，建议行胸部CT检查，以评估肺部炎症的范围和严重程度，及时发现病毒性肺炎或继发细菌性肺炎的影像学改变，为临床决策提供依据。

\subsection{治疗}

    甲型流感的治疗方式主要是药物治疗，需要进行一般治疗、对症治疗和抗病毒治疗。

    一般治疗主要通过卧床休息、多饮水、保持环境通风进行，对症治疗通常是针对咳嗽、流涕、发热等病症开具对应的综合性药物进行治疗。常见的治疗发热的药物是布洛芬，通过COX（环氧化酶）催化花生四烯酸转化为前列腺素，抑制 COX-1 和 COX-2 等因子的活性，减少 PGE2 合成，从而减缓炎症反应引起的发热。 

    下面主要讲解常用与抗甲流病毒的药物治疗途径。

    \subsubsection{神经氨酸酶抑制剂}

    流感病毒表面的血凝素 (HA) 和神经氨酸酶(NA) 在复制中各司其职，HA 主要负责粘附宿主细胞表面的唾液酸受体从而启动感染；待子代病毒在细胞表面组装完毕后，NA则负责切割细胞与病毒之间的乙酰神经氨酸，从而使子代流感病毒从宿主细胞表面释放。神经氨酸酶抑制剂（NAI）主要通过与神经氨酸酶的天然底物唾液酸竞争，选择性抑制病毒包膜上神经氨酸酶的活性，阻断酶活性位点，进而阻断病毒颗粒从被感染宿主细胞脱落，阻止病毒在宿主细胞间扩散，从而抑制病毒在体内的复制。

    WHO 和美国 CDC 推荐其为抗流感病毒一线治疗药物，联合使用抗病毒药物不能使病例临床获益。规范、足量、及早、足疗程使用时有助于避免流感病毒耐药病毒株的出现。
    
    注意：神经氨酸酶抑制剂禁止与非选择性胆碱酯酶抑制剂合用。

    \subsubsection{RNA 聚合酶抑制剂类药物}

    这类药物通过抑制病毒 RNA 聚合酶的活性，从根本上阻断病毒的复制。玛巴洛沙韦作用更为精准，能特异性抑制流感病毒 mRNA 合成所需的"帽状结构依赖性核酸内切酶"。

    这类药物常见的有玛巴洛沙韦和法维拉韦。他们的用药疗程都十分简单，不良反应较少或轻微，且玛巴洛沙韦在缓解症状方面与传统药物效果相当甚至更优。法维拉韦抗病毒谱广，对某些耐药病毒株都可能有效。

    \subsubsection{血凝素抑制剂}

    这类药物通过抑制病毒脂膜与宿主细胞的融合，来阻断病毒的复制过程。

\subsection{流感疫苗}

    \subsubsection{疫苗种类}

    疫苗是用无致病性或低致病性的病原体或其成分等制备的疫苗对宿主进行接种，诱导机体产生适应性免疫应答并获得长效的特异性免疫记忆的生物制剂。整体而言，疫苗包括第一代疫苗（传统疫苗）包括灭活疫苗、减毒活疫苗和类毒素；第二代疫苗（组分疫苗）包括由病原体的天然成分及其产物制成的亚单位疫苗、将能激发免疫应答的抗原成分的基因重组而产生的重组载体疫苗和结合疫苗；第三代疫苗（核酸疫苗）的代表为 DNA 疫苗、mRNA 疫苗。

    \begin{table}[H]
        \centering
        \caption{疫苗类型比较}
        \label{tab:vaccine_comparison}
        \small % 稍微缩小字体
        \begin{tabularx}{\textwidth}{@{}l >{\raggedright\arraybackslash}p{2.5cm} >{\raggedright\arraybackslash}p{2.8cm} >{\raggedright\arraybackslash}X >{\raggedright\arraybackslash}X@{}}
        \toprule
        \textbf{疫苗类型} & \textbf{制备方式} & \textbf{免疫应答类型} & \textbf{优点} & \textbf{缺点} \\
        \midrule
        灭活疫苗 & 病原体培养后经理化方法灭活 & 主要诱导特异性抗体（体液免疫） & 稳定、安全、易储存、制备简单快速 & 需多次接种；不能诱导CTL反应；可能引起局部或全身反应 \\
        \midrule
        减毒活疫苗 & 减毒或无毒的活病原体 & 体液免疫 + 细胞免疫 + 黏膜免疫（若经自然途径） & 免疫效果持久、良好 & 有回复突变风险（罕见）；免疫缺陷者和孕妇不宜接种 \\
        \midrule
        类毒素 & 外毒素经甲醛处理失活 & 诱导抗毒素（体液免疫） & 保留免疫原性，无毒 & 主要用于细菌外毒素相关疾病 \\
        \midrule
        亚单位疫苗 & 提取病原体有效免疫原成分（理化或DNA重组） & 体液免疫为主 & 安全、有效、成本低；不含活病原体或核酸 & 仅含部分抗原，可能免疫原性较弱 \\
        \midrule
        重组载体疫苗 & 编码免疫原基因插入减毒病毒/细菌载体 & 体液免疫 + 细胞免疫（强烈CTL反应） & 可诱导完整免疫应答；可设计多价疫苗 & 可能对载体本身产生免疫反应 \\
        \midrule
        结合疫苗 & 细菌荚膜多糖连接蛋白质载体 & 体液免疫（IgG类抗体） & 将TI抗原转为TD抗原，提高婴幼儿免疫效果 & 主要针对细菌荚膜多糖相关感染 \\
        \midrule
        DNA疫苗 & 编码免疫原基因与质粒构建重组体 & 体液免疫 + 细胞免疫 & 可持续表达、安全（无活病原体）、易制备 & 临床试验效果有限；仅适用于蛋白抗原 \\
        \midrule
        mRNA疫苗 & 编码免疫原mRNA + 脂质纳米颗粒包裹 & 体液免疫 + 细胞免疫（强佐剂效应） & 快速开发、强免疫应答、无整合风险、可组合多种抗原 & 不稳定，需低温储存 \\
        \bottomrule
        \end{tabularx}
    \end{table}

    具体到流感病毒，流感病毒疫苗有三种基本类型：
    \begin{itemize}
        \item 流感病毒灭活疫苗（IIV）
        \item 流感病毒减毒活疫苗（LAIV）
        \item 流感病毒重组疫苗（RIV）
    \end{itemize}
    
    其字母缩写后的数字表示疫苗效价（所含流感病毒血凝素 HA 抗原种类数），分为三价和四价。其中三价流感疫苗组份含有 A(H3N2) 亚型、A(H1N1)pdm09 亚型和 B 型流感病毒株的一个系，四价流感疫苗组份含 A(H3N2) 亚型、A(H1N1)pdm09 亚型和 B(Victoria) 系、B(Yamagata) 系。根据生产工艺，又可分为基于鸡胚培养、基于细胞培养和重组流感疫苗。根据世界卫生组织和美国疾病控制预防中心（CDC）的建议，2024–2025 年度推荐全面转向使用三价疫苗，这是因为既往季节曾推荐覆盖 B(Yamagata) 型病毒株的四价疫苗，然而，这种菌株已不再被认为在传播。三价重组流感疫苗（RIV3）和基于细胞培养的疫苗（ccIIV3）不含鸡蛋蛋白质。

    \subsubsection{抗原漂移与疫苗的有效时间}

    甲型流感病毒的 HA 和 NA 均容易发生变异，根据 HA 和 NA 的抗原性差异，可将甲型流感病毒分为若干亚型（subtype），目前 HA 包括 16 种亚型（H1\textasciitilde H16），NA 包括 9 种亚型（N1\textasciitilde N9）；

    甲型流感病毒的变异根据程度分为两种形式：ff 抗原漂移（antigenic drift）：抗原变异幅度小，主要是 HA 氨基酸的变异，其次是 NA 氨基酸的变异，属于量变。这种变异是由病毒基因点突变导致氨基酸变异，并可不断累积，由于人群存在一定的免疫保护作用，不会引起流感大流行，仅引起中、小规模流行，多出现在寒冷的季节，称为季节性流感。ff 抗原转换（antigenic shift）：指抗原变异幅度大，HA 或 NA 氨基酸的变异率达到 30\%\textasciitilde50\%，属于质变，常形成新的亚型（如 H1N1→H2N2）。这种变异多由基因重排形成的新亚型。如人群对流感病毒新亚型易感，并普遍缺乏针对新亚型的免疫力，可导致流感大流行。

    人体对感染流感病毒或接种流感疫苗后获得的免疫力会随时间衰减，衰减程度与人的年龄和身体状况、疫苗抗原等因素有关。澳大利亚开展的一项流感疫苗抗体动力学研究显示，接种流感疫苗 1 个月后各型疫苗株诱导的抗体水平达到高峰，3 个月左右后开始下降，在接种 6 个月后抗体水平仍高于基线，提示接种流感疫苗后抗体保护水平至少可维持 6 个月之久。另一方面，流行病学研究指出，流感疫苗触发的保护性免疫反应在几周内就会减弱。相对于其它病原体，随着人体内免疫记忆的丧失，流感疫苗失去有效性的速度更快。综合以上结果，指南推荐在每个流感活动高峰季节前 2-4 周安排接种工作，并对于不同年龄段的人群推荐了不同的接种种类和剂次。

    \begin{figure}[H]
        \centering
        \includegraphics[width=0.5\linewidth]{figures/Screenshot 2025-12-07 at 15.51.25.png}
        \caption{甲流疫苗有效性随着时间下降}
        \label{fig:vaccine-efficiency}
    \end{figure}

    \subsubsection{接种原则}

    \paragraph{建议优先接种人群}
    建议所有6月龄及以上且无接种禁忌的人都应接种流感疫苗。
    
    优先推荐重点和高风险人群包括：医务人员、60岁及以上的老年人、罹患一种或多种慢性病者（包括患有心血管疾病（单纯高血压除外）、慢性呼吸系统疾病、肝肾功能不全、血液病、神经系统疾病、神经肌肉功能障碍、代谢性疾病（包括糖尿病）等慢性病患者、患有免疫抑制疾病或免疫功能低下者）、养老机构、长期护理机构、福利院等聚集场所的人员、孕妇、6\textasciitilde59月龄的儿童、6月龄以下婴儿的看护人员、教学机构等监管场所的人员。
    
    \paragraph{接种时机}
    结合中国流感流行特点，9\textasciitilde10月一般是最佳接种时间，按照接种程序完成全程接种的人员，同一流感流行季节无需重复接种。
    
    \paragraph{接种禁忌与注意事项}
    对流感疫苗中所含任何成分（包括辅料、甲醛、裂解剂及抗生素）过敏者或有过任何一种流感疫苗接种严重过敏史者，禁止接种。
    
    以下人群禁止接种减毒活疫苗：
    
    \begin{itemize}
        \item 因使用药物、艾滋病毒感染 等任何原因造成免疫功能低下者；
        \item 长期使用含有阿司匹林或水杨酸成分药物治疗的儿童及青少年；
        \item 2\textasciitilde4岁患有哮喘的儿童；
        \item 妊娠女性；
        \item 有格林-巴利综合征病史者；
        \item 接种前48小时使用过奥司他韦、扎那米韦等抗病毒药物者，或接种前5天使用过帕拉米韦，或接种前17天使用过巴洛沙韦者。
    \end{itemize}
    
    流感病毒灭活疫苗与其他灭活疫 苗或减毒活疫苗可同时在不同部位接种。但接种流感减毒活疫苗后必须间隔4周以上才可接种其他减毒活疫苗。
    
    患有急性疾病、严重慢性疾病或慢性疾病的急性 发作期以及发热患者，建议痊愈或者病情稳定控制后接种。
    
    对于案例患者为60岁以上老年人，且患有慢性病，因而属于优先接种人群，建议接种流感疫苗。


    \subsection{重症流感及并发症}

    \subsubsection{重症流感的定义}

    成人流感满足以下任何一条标准之一即为重症流感：
    \begin{itemize}
        \item 呼吸急促，$\mathrm{RR}\ge30 \text{次/分}$
        \item 静息状态下，吸空气时指$\text{氧饱和度}\le93$\%
        \item 动脉血氧分压（$\text{PaO}_2$）/吸氧浓度（$\text{FiO}_2$）$\le300$，高海拔（海拔超过 1000 米）地区应根据以下公式对 $\text{PaO}_2/\text{FiO}_2$ 进行校正：PaO$_2$/FiO$_2$×[760/大气压（mmHg）]（1mmHg=0.133kPa）
        \item 临床症状进行性加重，肺部影像学显示 24\textasciitilde48 小时内病灶明显进展＞50\%
    \end{itemize}

    儿童符合下列任何一条：

    \begin{itemize}
        \item 超高热或持续高热超过 3 天
        \item 呼吸急促（＜2 月龄，RR$\ge60$ 次/分；2\textasciitilde12 月龄，RR$\ge$50 次/分；1\textasciitilde5 岁，RR$\ge$40 次/分；＞5 岁，RR$\ge$30 次/分），除外发热和哭闹的影响
        \item 静息状态下，吸空气时指氧饱和度$\le$93\%
        \item 鼻翼扇动、三凹征、喘鸣或喘息
        \item 意识障碍或惊厥
        \item 拒食或喂养困难，有脱水征
    \end{itemize}
    
    有下列情况之一者为危重症：
    \begin{itemize}
        \item 呼吸衰竭，且需要机械通气
        \item 休克
        \item 急性坏死性脑病
        \item 合并其他器官功能衰竭需 ICU 监护治疗
    \end{itemize}

    \subsubsection{常见的并发症}

    流行性感冒最常见的并发症即病毒性肺炎，其他并发症包括：

    \paragraph{继发性细菌性肺炎}
    流感病毒可侵犯下呼吸道，引起原发性病毒性肺炎。部分流感患者可合并细菌、其他病毒、非典型病原体、真菌等感染。合并金黄色葡萄球菌、肺炎链球菌或侵袭性肺曲霉感染时，病情重、病死率高。

    \paragraph{神经系统并发症}
    神经系统损伤包括脑炎、脑病、脊髓炎、吉兰-巴雷综合征（Guillain-Barre Syndrome）等，儿童多于成人，急性坏死性脑病更为凶险。

    \paragraph{心血管系统并发症}
    心脏损伤主要有心肌炎、心包炎。可见心肌标志物、心电图、心脏超声等异常，严重者可出现心力衰竭。此外，感染流感病毒后，心肌梗死、缺血性心脏病相关住院和死亡的风险明显增加。

    \paragraph{肌肉与免疫系统并发症}
    肌炎和横纹肌溶解主要表现为肌痛、肌无力、血清肌酸激酶、肌红蛋白升高，严重者可导致急性肾损伤等。

    \subsection{流感病毒的传播}

    长期以来，公共卫生指南普遍认为流感病毒主要通过两大途径传播：
    \begin{itemize}
    \item \textbf{飞沫传播（Droplet Transmission）}：感染者在咳嗽、打喷嚏或说话时，会产生含有病毒的大颗粒呼吸道飞沫（通常直径 >5 微米）。这些飞沫在短距离内（通常为 1-2 米）直接沉降到易感者的口、鼻或眼部黏膜上，导致感染。基于此理论，保持社交距离和佩戴医用口罩是核心预防措施。
    \item \textbf{接触传播（Contact Transmission）}：含有病毒的呼吸道飞沫污染物体表面（如门把手、桌面），健康人接触这些被污染的表面后，再触摸自己的口、鼻、眼，从而发生间接感染。因此，频繁和正确的手部卫生被视为切断接触传播链的关键措施。
    \end{itemize}

    这些传统途径在多份历史性的感染控制指南中被反复强调，并构成了标准预防（Standard Precautions）和飞沫预防（Droplet Precautions）措施的核心依据。

    \textbf{日益受到重视的传播途径：气溶胶传播}

    近年来，越来越多的证据表明，气溶胶传播（Aerosol Transmission）在流感病毒的传播中扮演着比以往认知中更为重要的角色。

    \begin{itemize}
    \item \textbf{气溶胶的定义与产生}：气溶胶是指直径小于 5 微米的微小呼吸道颗粒，由感染者通过呼吸、说话、咳嗽乃至歌唱等活动产生。与大飞沫不同，气溶胶可以在空气中悬浮数小时，并能传播到更远的距离，尤其是在通风不良的室内环境中。
    \item \textbf{支持气溶胶传播的证据}：多项研究，包括对医院和学校流感爆发的调查，均发现了支持气溶胶传播的证据。研究指出，不利的通风条件，例如通风系统关闭或空气流动不畅，可能会显著促进病毒的气溶胶传播。这一传播模式的存在意味着，仅依赖短距离社交隔离和常规医用口罩可能不足以完全阻断传播，尤其是在人员密集的室内空间。
    \item \textbf{对防控策略的启示}：对气溶胶传播的认可，极大地推动了对环境干预措施的重视，尤其是改善室内通风和空气过滤。此外，这也为在特定高风险情境下（如对患者进行气溶胶生成操作时）推荐使用 N95 或同等级别的呼吸防护具提供了理论基础。尽管关于不同途径的相对贡献仍在争议中，但将气溶胶传播纳入综合防控考量已成为学界共识。
    \end{itemize}

\section{贡献}

    这份报告由作者根据2025秋季学期 PBL-G1 课程班级各个小组的作业整理而成。

    参与贡献的班级成员：
    \begin{itemize}
        \item \namebox{\textbf{蒋子睿}}——\textit{临床四班}
        \item \namebox{\textbf{王\hfill 涵}}——\textit{临床四班}
        \item \namebox{\textbf{罗\hfill 昊}}——\textit{临床四班}
        \item \namebox{\textbf{金洪民}}——\textit{临床六班}
        \item \namebox{\textbf{胡梦玲}}——\textit{临床六班}
        \item \namebox{\textbf{林雨朵}}——\textit{临床七班}
        \item \namebox{\textbf{周钲珺}}——\textit{临床七班}
        \item \namebox{\textbf{张元畅}}——\textit{临床九班}
    \end{itemize}

    指导老师：\namebox{\textbf{朱翔}}
%----------------------------------------------------------

% 加入引用
% \printbibliography

%----------------------------------------------------------

\end{document}